% Methods appendix for the TMLR version (prose; the tables are in appendix_tables.tex).

\section{Task and data}\label{app:data}
Each scene is a chain of three objects, R, M and C, on an integer grid bounded at $\pm 6$. M is placed relative to R
and C relative to M, each by one to five units along one axis (with probability 0.7) or along both axes (0.3), and no
two objects share a position. Each placement is written with one of four templates (``X is d of Y'', ``X lies d of
Y'', ``X sits d of Y'', ``X is located d of Y''), where d is an offset such as ``4 units east'' or ``2 units north and
1 unit west''. Names are single tokens in the Gemma 3 tokenizer, drawn from a pool of person names and a pool of object
names that alternate between scenes. The description is followed by the object list in random order (``The objects
are: \dots'') and by one question, ``Which is farther east, X or Y?'' or ``Which is farther north, X or Y?'', with C and
R in both orders, then ``Answer with one word.'' The whole text is wrapped in the model's chat template. We keep only
questions in which C and R differ by at least one unit on the asked axis, and fill each kind of question (own step only,
agree, first step only, conflict with C's step larger, conflict with M's step larger) up to its preregistered size. The
model's answer is the sign of the difference between the logits of the two names, which are single tokens, at the
final position.

\section{Probes}\label{app:probes}
Probes are ridge regressions on the residual stream, fitted separately on east and on north questions. The scenes are
split once into 80\% for training and 20\% for testing (fixed seed), so that every question of a test scene is
unseen, and the ridge penalty is chosen by cross-validation grouped by scene within the training scenes. The step
probes predict $d_1$ and $d_2$ and are scored by held-out $R^2$, averaged over the two axes. The truth probe predicts
the correct answer coded as $\pm 1$ and is scored by the sign of its prediction. The probed sites are the last token of
M's sentence and of C's sentence, C's and R's names in the question, the token after the second of those names, and the
final position, at every second layer (Tables~\ref{tab:app-flow12-r2} and~\ref{tab:app-flow27-r2}). In the
composition test (Table~\ref{tab:app-compose}) the target is the distance between a two-step object and the first
object, and the additive baseline is a ridge regression on a one-hot code of each of the two sentences' offsets.

\section{Activation patching}\label{app:patching}
For each kind of question and each of the two steps, 100 questions are paired with a counterfactual in which that step
is negated on the asked axis; objects placed relative to the moved one move with it, and the words are otherwise
unchanged. We copy the residual stream of the counterfactual run into the clean run at one site and one layer (every
fourth layer and the last but two) and report the change in the logit difference as a share of the full
counterfactual change (Tables~\ref{tab:app-copy12own}--\ref{tab:app-copy27first}). The four sites are all tokens of
M's sentence, all tokens of C's sentence, all tokens between the object list and the final position (the question,
``Answer with one word.'' and the chat-template tokens), and the final position.

\section{Direct attribution and interventions}\label{app:attribution}
Direct attribution writes the final logit difference as a sum over the embeddings, each attention head (its slice of
the attention output passed through the output projection) and each MLP, with every normalisation scale taken from the
run itself, including Gemma's post-attention and post-MLP norms and the final norm. Each head's contribution is further
split by source token using its attention weights and value vectors, and the sources are grouped into M's sentence, C's
sentence, the object list, C's and R's names in the question, the other question and template tokens, and the final
position. On the discovery set each component's contribution is regressed on own\_yes, $d_1$, $d_2$ and the name order,
and the components are ranked by their own\_yes coefficient (Tables~\ref{tab:app-writers12}
and~\ref{tab:app-writers27}). The interventions act on the same quantities at the final position only: \emph{mean
ablation} replaces a component's output by its mean over questions, \emph{twin-mean removal} by the average of its
outputs on the question and on its twin, and a \emph{swap} by its output on the twin. Matched random controls draw
components of the same kind from the same layers (MLPs within four layers), excluding the top writers, in three draws.
To see where the removed shortcut comes back, we repeat the attribution with the intervention in place on 600 questions
of the removal set, measure every component at the clean run's final-norm scale, and count the rescaling by the final
norm separately (Tables~\ref{tab:app-repair12} and~\ref{tab:app-repair27}).

\section{Compute}\label{app:compute}
All runs used one GPU on Google Colab (A100 40 GB for 12B, H100 or 80 GB-class GPUs for 27B, and an L4 with part of
the 12B model in CPU memory for the order-and-voice control) with bfloat16 weights; attribution used the eager
attention implementation so that attention weights are available. Each test took between about fifteen minutes and a
few hours.

\section{Result tables}\label{app:tables}
The tables below are generated directly from the project's result log, so that every number matches the released
result files.

% Generated by make_appendix.py from results/RESULTS.md; do not edit by hand.

\begin{table}[htbp]
\centering
\scriptsize
\caption{Behaviour on the 5,250-question set (Gemma 3 12B), clean and under two earlier preregistered erasures: the probe-readable coordinate directions (layers 12--48, every token but the first) and a variance-matched random erasure. Excess drop = random-erasure accuracy minus map-erased accuracy; 95\% scene-bootstrap CIs.}
\label{tab:app-bulk}
\begin{tabular}{@{}lrrrrr@{}}
\toprule
cell & n & clean & map erased & fair random & \begin{tabular}[b]{@{}c@{}}excess drop\\{}[95\% CI]\end{tabular} \\
\midrule
own step only & 750 & 0.98 & 0.89 & 0.94 & +0.048 [+0.023, +0.075] \\
agree & 750 & 0.97 & 0.91 & 0.94 & +0.027 [+0.007, +0.047] \\
first step only & 750 & 0.53 & 0.52 & 0.57 & +0.057 [+0.043, +0.075] \\
conflict, C larger & 1500 & 0.88 & 0.68 & 0.83 & +0.147 [+0.117, +0.175] \\
conflict, M larger \textbf{(decisive)} & 1500 & 0.34 & 0.45 & 0.38 & -0.071 [-0.093, -0.049] \\
all (balanced pool) & 5250 & 0.70 & 0.65 & 0.70 & +0.041 [+0.029, +0.052] \\
\bottomrule
\end{tabular}
\end{table}

\begin{table}[htbp]
\centering
\scriptsize
\caption{Composition test (12B; 3,000 scenes of 3--6 objects, no question). Probes for each two-step object's distance and absolute coordinates relative to the first object, at its mention in the object list and at the end of its own sentence, against the best additive code of the two sentences (one-hot offsets).}
\label{tab:app-compose}
\begin{tabular}{@{}lrrrrr@{}}
\toprule
site & layer & target & probe R² & \begin{tabular}[b]{@{}c@{}}parts\\{}baseline R²\end{tabular} & gap [95\% CI] \\
\midrule
list mention & 21 & dist & 0.29 & 0.15 & +0.14 [+0.07, +0.22] \\
list mention & 21 & abs\_x & 0.23 & 0.21 & +0.02 [-0.05, +0.09] \\
list mention & 21 & abs\_y & 0.24 & 0.21 & +0.03 [-0.05, +0.11] \\
sentence end & 24 & dist & 0.61 & 0.15 & +0.45 [+0.38, +0.53] \\
sentence end & 24 & abs\_x & 0.53 & 0.21 & +0.32 [+0.25, +0.40] \\
sentence end & 24 & abs\_y & 0.56 & 0.21 & +0.35 [+0.27, +0.43] \\
\bottomrule
\end{tabular}
\end{table}

\begin{table}[htbp]
\centering
\scriptsize
\caption{Flow set, Gemma 3 12B: held-out $R^2$ of M's step $d_1$ / C's step $d_2$ by site and layer.}
\label{tab:app-flow12-r2}
\begin{tabular}{@{}lrrrrrr@{}}
\toprule
layer & M\_end & C\_end & Q\_C & Q\_R & Q\_after & final \\
\midrule
0 & -0.02 / -0.02 & -0.02 / -0.02 & -0.02 / -0.02 & -0.02 / -0.01 & -0.02 / -0.02 & -0.02 / -0.02 \\
4 & 0.72 / -0.03 & 0.55 / 0.45 & -0.02 / -0.02 & -0.02 / -0.01 & 0.08 / -0.03 & 0.28 / -0.06 \\
8 & 0.73 / 0.02 & 0.52 / 0.55 & -0.02 / -0.02 & 0.09 / -0.01 & 0.55 / 0.42 & 0.52 / -0.02 \\
12 & 0.72 / 0.01 & 0.49 / 0.63 & 0.09 / 0.16 & 0.36 / 0.27 & 0.66 / 0.60 & 0.57 / 0.40 \\
16 & 0.78 / 0.03 & 0.71 / 0.74 & 0.68 / 0.67 & 0.70 / 0.69 & 0.80 / 0.82 & 0.76 / 0.71 \\
20 & 0.79 / 0.03 & 0.75 / 0.82 & 0.75 / 0.80 & 0.75 / 0.76 & 0.84 / 0.85 & 0.81 / 0.82 \\
24 & 0.77 / -0.00 & 0.75 / 0.82 & 0.79 / 0.87 & 0.80 / 0.83 & 0.82 / 0.86 & 0.82 / 0.86 \\
28 & 0.75 / -0.01 & 0.69 / 0.77 & 0.80 / 0.86 & 0.79 / 0.81 & 0.83 / 0.87 & 0.84 / 0.89 \\
32 & 0.74 / -0.00 & 0.61 / 0.71 & 0.76 / 0.84 & 0.76 / 0.76 & 0.80 / 0.85 & 0.81 / 0.85 \\
36 & 0.68 / -0.00 & 0.44 / 0.61 & 0.73 / 0.80 & 0.73 / 0.73 & 0.77 / 0.81 & 0.76 / 0.81 \\
40 & 0.64 / -0.01 & 0.39 / 0.58 & 0.73 / 0.79 & 0.71 / 0.72 & 0.75 / 0.80 & 0.74 / 0.80 \\
44 & 0.60 / -0.02 & 0.27 / 0.53 & 0.73 / 0.78 & 0.71 / 0.72 & 0.71 / 0.76 & 0.67 / 0.74 \\
48 & 0.66 / -0.02 & 0.16 / 0.54 & 0.69 / 0.75 & 0.69 / 0.69 & 0.65 / 0.69 & 0.64 / 0.70 \\
\bottomrule
\end{tabular}
\end{table}

\begin{table}[htbp]
\centering
\scriptsize
\caption{Flow set, Gemma 3 27B: held-out $R^2$ of $d_1$ / $d_2$ by site and layer.}
\label{tab:app-flow27-r2}
\begin{tabular}{@{}lrrrrrr@{}}
\toprule
layer & M\_end & C\_end & Q\_C & Q\_R & Q\_after & final \\
\midrule
0 & -0.02 / -0.02 & -0.02 / -0.02 & -0.02 / -0.02 & -0.02 / -0.01 & -0.02 / -0.02 & -0.02 / -0.02 \\
4 & 0.71 / -0.02 & 0.36 / 0.53 & -0.02 / -0.01 & -0.02 / -0.01 & 0.24 / -0.02 & 0.11 / 0.09 \\
8 & 0.74 / -0.03 & 0.44 / 0.60 & -0.02 / -0.02 & -0.02 / -0.01 & 0.28 / 0.06 & 0.33 / 0.08 \\
12 & 0.75 / 0.02 & 0.60 / 0.67 & -0.05 / -0.01 & -0.02 / -0.01 & 0.52 / 0.36 & 0.47 / 0.28 \\
16 & 0.71 / 0.03 & 0.54 / 0.70 & 0.03 / -0.01 & -0.02 / -0.01 & 0.67 / 0.56 & 0.62 / 0.40 \\
20 & 0.78 / 0.04 & 0.63 / 0.75 & 0.57 / 0.39 & 0.60 / 0.44 & 0.80 / 0.69 & 0.69 / 0.51 \\
24 & 0.81 / 0.04 & 0.72 / 0.82 & 0.75 / 0.79 & 0.78 / 0.75 & 0.86 / 0.89 & 0.79 / 0.74 \\
28 & 0.82 / 0.05 & 0.78 / 0.85 & 0.84 / 0.89 & 0.84 / 0.84 & 0.89 / 0.90 & 0.87 / 0.86 \\
32 & 0.83 / 0.02 & 0.78 / 0.85 & 0.86 / 0.91 & 0.85 / 0.86 & 0.89 / 0.90 & 0.88 / 0.92 \\
36 & 0.80 / 0.03 & 0.77 / 0.84 & 0.85 / 0.91 & 0.84 / 0.84 & 0.88 / 0.92 & 0.89 / 0.92 \\
40 & 0.77 / 0.03 & 0.74 / 0.80 & 0.83 / 0.89 & 0.82 / 0.82 & 0.88 / 0.92 & 0.88 / 0.92 \\
44 & 0.74 / 0.01 & 0.69 / 0.77 & 0.82 / 0.89 & 0.81 / 0.81 & 0.87 / 0.92 & 0.85 / 0.91 \\
48 & 0.73 / -0.01 & 0.54 / 0.68 & 0.79 / 0.87 & 0.79 / 0.78 & 0.85 / 0.89 & 0.84 / 0.89 \\
52 & 0.70 / -0.01 & 0.54 / 0.67 & 0.76 / 0.87 & 0.79 / 0.77 & 0.83 / 0.89 & 0.84 / 0.89 \\
56 & 0.66 / -0.02 & 0.49 / 0.62 & 0.76 / 0.85 & 0.76 / 0.74 & 0.82 / 0.88 & 0.81 / 0.86 \\
60 & 0.57 / -0.02 & 0.46 / 0.61 & 0.76 / 0.84 & 0.75 / 0.73 & 0.78 / 0.85 & 0.79 / 0.84 \\
\bottomrule
\end{tabular}
\end{table}

\begin{table}[htbp]
\centering
\scriptsize
\caption{Logit lens, 12B: share of questions on which the layer's C-versus-R logit points to the correct answer.}
\label{tab:app-lens12}
\begin{tabular}{@{}lrrrrr@{}}
\toprule
layer & own step only & agree & \begin{tabular}[b]{@{}c@{}}first\\{}step only\end{tabular} & \begin{tabular}[b]{@{}c@{}}conflict,\\{}C larger\end{tabular} & \begin{tabular}[b]{@{}c@{}}conflict,\\{}M larger\end{tabular} \\
\midrule
24 & 0.55 & 0.48 & 0.47 & 0.56 & 0.56 \\
28 & 0.56 & 0.53 & 0.45 & 0.52 & 0.52 \\
32 & 0.65 & 0.66 & 0.51 & 0.58 & 0.48 \\
36 & 0.98 & 0.96 & 0.51 & 0.86 & 0.39 \\
40 & 0.98 & 0.96 & 0.50 & 0.87 & 0.38 \\
44 & 0.97 & 0.97 & 0.51 & 0.87 & 0.40 \\
48 & 0.99 & 0.97 & 0.50 & 0.89 & 0.40 \\
\bottomrule
\end{tabular}
\end{table}

\begin{table}[htbp]
\centering
\scriptsize
\caption{Logit lens, 27B.}
\label{tab:app-lens27}
\begin{tabular}{@{}lrrrrr@{}}
\toprule
layer & own step only & agree & \begin{tabular}[b]{@{}c@{}}first\\{}step only\end{tabular} & \begin{tabular}[b]{@{}c@{}}conflict,\\{}C larger\end{tabular} & \begin{tabular}[b]{@{}c@{}}conflict,\\{}M larger\end{tabular} \\
\midrule
32 & 0.54 & 0.60 & 0.49 & 0.53 & 0.51 \\
36 & 0.54 & 0.58 & 0.55 & 0.54 & 0.49 \\
40 & 0.57 & 0.59 & 0.54 & 0.55 & 0.50 \\
44 & 0.90 & 0.85 & 0.57 & 0.82 & 0.38 \\
48 & 0.97 & 0.89 & 0.54 & 0.90 & 0.31 \\
52 & 0.97 & 0.87 & 0.57 & 0.90 & 0.32 \\
56 & 0.99 & 0.90 & 0.56 & 0.95 & 0.27 \\
60 & 0.99 & 0.90 & 0.55 & 0.96 & 0.23 \\
\bottomrule
\end{tabular}
\end{table}

\begin{table}[htbp]
\centering
\scriptsize
\caption{Flow set, 12B: change in the logit difference when one step is flipped (oriented so that + is the way the flipped step points).}
\label{tab:app-flip12}
\begin{tabular}{@{}lrrr@{}}
\toprule
flipped step & cell & n & full effect \\
\midrule
own & own step only & 100 & +31.10 [+29.44, +32.80] \\
own & agree & 98 & +26.77 [+23.57, +30.08] \\
own & conflict, C larger & 100 & +29.97 [+27.24, +32.32] \\
own & conflict, M larger & 100 & +20.55 [+17.37, +23.99] \\
parent & agree & 98 & +9.76 [+7.25, +12.43] \\
parent & first step only & 100 & +2.64 [+1.65, +3.76] \\
parent & conflict, C larger & 100 & +5.90 [+3.80, +8.30] \\
parent & conflict, M larger & 100 & +12.11 [+9.16, +15.07] \\
\bottomrule
\end{tabular}
\end{table}

\begin{table}[htbp]
\centering
\scriptsize
\caption{Flow set, 27B: as Table~\ref{tab:app-flip12}.}
\label{tab:app-flip27}
\begin{tabular}{@{}lrrr@{}}
\toprule
flipped step & cell & n & full effect \\
\midrule
own & own step only & 100 & +58.60 [+56.18, +61.18] \\
own & agree & 98 & +39.85 [+34.72, +44.95] \\
own & conflict, C larger & 100 & +49.81 [+46.46, +52.71] \\
own & conflict, M larger & 100 & +33.02 [+29.40, +36.97] \\
parent & agree & 98 & +8.68 [+5.87, +11.81] \\
parent & first step only & 100 & +7.82 [+6.13, +9.55] \\
parent & conflict, C larger & 100 & +3.59 [+1.47, +6.01] \\
parent & conflict, M larger & 100 & +12.38 [+8.92, +15.98] \\
\bottomrule
\end{tabular}
\end{table}

\begin{table}[htbp]
\centering
\scriptsize
\caption{12B: share of the effect of flipping C's own step carried by copying one site at one layer.}
\label{tab:app-copy12own}
\begin{tabular}{@{}lrrrr@{}}
\toprule
layer & M\_sentence & C\_sentence & question & final \\
\midrule
0 & 0.00 & 1.00 & 0.00 & 0.00 \\
4 & 0.00 & 1.00 & -0.00 & 0.00 \\
8 & 0.00 & 1.00 & 0.01 & 0.00 \\
12 & 0.00 & 1.00 & 0.02 & 0.00 \\
16 & 0.00 & 0.91 & 0.11 & 0.00 \\
20 & 0.00 & 0.26 & 0.74 & 0.00 \\
24 & 0.00 & -0.00 & 0.96 & 0.04 \\
28 & 0.00 & -0.00 & 0.49 & 0.51 \\
32 & 0.00 & 0.00 & 0.29 & 0.71 \\
36 & 0.00 & 0.00 & 0.20 & 0.80 \\
40 & 0.00 & 0.00 & 0.20 & 0.80 \\
44 & 0.00 & 0.00 & 0.19 & 0.82 \\
46 & 0.00 & 0.00 & 0.08 & 0.92 \\
\bottomrule
\end{tabular}
\end{table}

\begin{table}[htbp]
\centering
\scriptsize
\caption{12B: share of the effect of flipping M's step carried by copying one site at one layer.}
\label{tab:app-copy12first}
\begin{tabular}{@{}lrrrr@{}}
\toprule
layer & M\_sentence & C\_sentence & question & final \\
\midrule
0 & 1.00 & -0.00 & -0.00 & -0.00 \\
4 & 1.01 & -0.00 & -0.01 & -0.00 \\
8 & 0.98 & 0.04 & 0.01 & -0.00 \\
12 & 0.92 & 0.20 & -0.01 & -0.00 \\
16 & 0.67 & 0.45 & 0.10 & -0.00 \\
20 & 0.22 & 0.29 & 0.53 & -0.00 \\
24 & -0.01 & 0.01 & 0.95 & 0.03 \\
28 & -0.00 & -0.01 & 0.53 & 0.48 \\
32 & 0.00 & -0.01 & 0.33 & 0.67 \\
36 & 0.01 & -0.01 & 0.24 & 0.75 \\
40 & 0.01 & -0.01 & 0.24 & 0.76 \\
44 & 0.00 & -0.00 & 0.22 & 0.78 \\
46 & 0.00 & -0.00 & 0.09 & 0.92 \\
\bottomrule
\end{tabular}
\end{table}

\begin{table}[htbp]
\centering
\scriptsize
\caption{27B: as Table~\ref{tab:app-copy12own}.}
\label{tab:app-copy27own}
\begin{tabular}{@{}lrrrr@{}}
\toprule
layer & M\_sentence & C\_sentence & question & final \\
\midrule
0 & 0.00 & 1.00 & 0.00 & 0.00 \\
4 & -0.00 & 1.00 & -0.00 & 0.00 \\
8 & -0.00 & 1.00 & 0.00 & 0.00 \\
12 & -0.00 & 1.00 & 0.00 & 0.00 \\
16 & -0.00 & 1.00 & 0.01 & 0.00 \\
20 & 0.00 & 0.98 & 0.05 & 0.00 \\
24 & -0.00 & 0.61 & 0.44 & 0.00 \\
28 & 0.00 & 0.31 & 0.70 & 0.01 \\
32 & 0.00 & 0.02 & 0.72 & 0.30 \\
36 & 0.00 & 0.00 & 0.29 & 0.74 \\
40 & 0.00 & 0.00 & 0.14 & 0.85 \\
44 & 0.00 & 0.00 & 0.13 & 0.87 \\
48 & 0.00 & -0.00 & 0.12 & 0.87 \\
52 & 0.00 & 0.00 & 0.12 & 0.87 \\
56 & 0.00 & 0.00 & 0.09 & 0.91 \\
60 & 0.00 & 0.00 & 0.04 & 0.96 \\
\bottomrule
\end{tabular}
\end{table}

\begin{table}[htbp]
\centering
\scriptsize
\caption{27B: as Table~\ref{tab:app-copy12first}.}
\label{tab:app-copy27first}
\begin{tabular}{@{}lrrrr@{}}
\toprule
layer & M\_sentence & C\_sentence & question & final \\
\midrule
0 & 1.00 & -0.00 & -0.00 & -0.00 \\
4 & 1.00 & -0.00 & -0.00 & -0.00 \\
8 & 1.01 & -0.01 & -0.02 & -0.00 \\
12 & 1.02 & -0.01 & -0.01 & -0.00 \\
16 & 0.94 & 0.07 & 0.01 & -0.00 \\
20 & 0.86 & 0.20 & -0.04 & -0.00 \\
24 & 0.24 & 0.29 & 0.21 & -0.00 \\
28 & 0.10 & 0.24 & 0.57 & 0.03 \\
32 & 0.00 & 0.02 & 0.66 & 0.32 \\
36 & 0.00 & 0.00 & 0.42 & 0.57 \\
40 & 0.00 & 0.00 & 0.25 & 0.68 \\
44 & 0.01 & 0.00 & 0.21 & 0.73 \\
48 & 0.01 & -0.01 & 0.19 & 0.76 \\
52 & 0.01 & -0.00 & 0.19 & 0.77 \\
56 & -0.00 & -0.00 & 0.13 & 0.84 \\
60 & -0.00 & -0.00 & 0.04 & 0.96 \\
\bottomrule
\end{tabular}
\end{table}

\begin{table}[htbp]
\centering
\scriptsize
\caption{12B, discovery set: the 20 largest writers of the shortcut (own\_yes coefficient of each component's direct contribution), their weights on $d_1$ and $d_2$, and the token groups their shortcut comes from.}
\label{tab:app-writers12}
\begin{tabular}{@{}lrrrrr@{}}
\toprule
component & own\_yes & d1 & d2 & mean & \begin{tabular}[b]{@{}c@{}}own\_yes\\{}comes from\end{tabular} \\
\midrule
L44H10 & +6.61 & +0.357 & +0.195 & -1.69 & Q\_other 0.80, final 0.13 \\
L47H2 & +4.64 & +0.213 & +0.125 & -1.40 & Q\_C 0.30, list 0.27 \\
L45H1 & +4.47 & +0.269 & +0.199 & -1.43 & Q\_other 0.79, final 0.13 \\
L41H0 & +3.73 & +0.540 & +0.262 & -1.57 & Q\_C 0.35, list 0.33 \\
L41H11 & +3.48 & +0.184 & +0.112 & -0.88 & list 0.29, C\_sentence 0.28 \\
L42H0 & +2.76 & +0.130 & +0.090 & -0.66 & Q\_other 0.83, final 0.08 \\
L41H5 & +2.60 & +0.038 & +0.243 & -1.32 & list 0.38, Q\_C 0.28 \\
L41H3 & +2.60 & +0.096 & +0.089 & -0.77 & list 0.31, Q\_C 0.26 \\
L35H8 & +2.50 & +0.127 & +0.035 & -0.61 & Q\_C 0.37, list 0.23 \\
L47H10 & +2.47 & +0.130 & +0.091 & -0.90 & Q\_C 0.31, list 0.25 \\
L47H8 & +1.95 & +0.103 & +0.062 & -0.41 & Q\_other 0.52, final 0.48 \\
L35H3 & +1.66 & +0.080 & +0.079 & -0.55 & list 0.41, Q\_C 0.25 \\
L41H4 & +1.61 & +0.045 & -0.134 & -0.14 & Q\_C 0.35, list 0.30 \\
MLP40 & +1.48 & +0.085 & +0.061 & -0.48 & – \\
L35H12 & +1.45 & +0.050 & +0.008 & -0.35 & Q\_C 0.36, list 0.28 \\
L46H4 & +1.43 & +0.049 & +0.061 & -0.51 & – \\
L47H7 & +1.39 & +0.060 & +0.039 & -0.36 & – \\
MLP43 & +1.20 & +0.047 & -0.075 & -0.29 & – \\
L41H1 & +1.13 & +0.050 & +0.079 & -0.11 & – \\
L32H7 & +1.12 & +0.050 & +0.036 & -0.32 & – \\
\bottomrule
\end{tabular}
\end{table}

\begin{table}[htbp]
\centering
\scriptsize
\caption{27B: as Table~\ref{tab:app-writers12}.}
\label{tab:app-writers27}
\begin{tabular}{@{}lrrrrr@{}}
\toprule
component & own\_yes & d1 & d2 & mean & \begin{tabular}[b]{@{}c@{}}own\_yes\\{}comes from\end{tabular} \\
\midrule
L57H18 & +7.26 & +0.311 & +0.155 & -1.08 & Q\_other 0.75, final 0.18 \\
L59H28 & +5.95 & +0.194 & +0.266 & -0.92 & list 0.23, Q\_C 0.23 \\
L59H27 & +3.47 & +0.128 & +0.186 & -0.50 & list 0.23, Q\_C 0.20 \\
L47H12 & +2.96 & +0.110 & +0.141 & -0.47 & Q\_C 0.34, list 0.24 \\
L58H19 & +2.66 & +0.071 & +0.171 & -0.26 & Q\_other 0.73, final 0.24 \\
L53H11 & +2.60 & +0.022 & +0.006 & -0.35 & list 0.27, Q\_C 0.23 \\
L60H30 & +2.42 & +0.095 & +0.111 & -0.30 & final 0.68, Q\_other 0.32 \\
L55H28 & +2.41 & +0.170 & +0.212 & -0.34 & list 0.34, Q\_C 0.21 \\
L53H27 & +2.18 & +0.167 & +0.157 & -0.78 & list 0.28, Q\_R 0.25 \\
L58H15 & +1.94 & +0.059 & +0.093 & -0.38 & list 0.34, C\_sentence 0.21 \\
MLP52 & +1.91 & +0.090 & +0.018 & -0.36 & – \\
L54H0 & +1.89 & +0.050 & +0.047 & -0.29 & Q\_other 0.78, final 0.14 \\
MLP54 & +1.88 & +0.124 & +0.109 & -0.30 & – \\
L53H30 & +1.85 & +0.061 & +0.050 & -0.14 & list 0.29, Q\_C 0.26 \\
L59H4 & +1.71 & +0.043 & +0.087 & -0.38 & Q\_C 0.25, list 0.22 \\
L47H22 & +1.56 & +0.023 & +0.121 & -0.37 & – \\
L56H1 & +1.56 & +0.032 & +0.052 & -0.33 & – \\
L55H9 & +1.51 & +0.060 & -0.067 & -0.20 & – \\
L47H26 & +1.48 & +0.003 & -0.069 & -0.24 & – \\
L58H9 & +1.48 & -0.013 & +0.015 & -0.14 & – \\
\bottomrule
\end{tabular}
\end{table}

\begin{table}[htbp]
\centering
\scriptsize
\caption{12B, writers' test set: mean ablation at the final position (accuracy and regression own\_yes / $d_1$ / $d_2$).}
\label{tab:app-meanabl12}
\begin{tabular}{@{}lrrrrrr@{}}
\toprule
condition & own step only & agree & \begin{tabular}[b]{@{}c@{}}first\\{}step only\end{tabular} & \begin{tabular}[b]{@{}c@{}}conflict,\\{}C larger\end{tabular} & \begin{tabular}[b]{@{}c@{}}conflict,\\{}M larger\end{tabular} & \begin{tabular}[b]{@{}c@{}}regression:\\{}own\_yes / d1 / d2\end{tabular} \\
\midrule
clean & 0.98 & 0.98 & 0.47 & 0.92 & 0.34 & +19.8 / +1.25 / +1.37 \\
own5 & 0.95 & 0.98 & 0.47 & 0.86 & 0.34 & +19.2 / +1.19 / +1.40 \\
own10 & 0.95 & 0.96 & 0.47 & 0.85 & 0.35 & +21.3 / +1.30 / +1.61 \\
own20 & 0.89 & 0.91 & 0.47 & 0.82 & 0.34 & +12.4 / +0.63 / +0.87 \\
random & 0.98 & 0.98 & 0.47 & 0.90 & 0.34 & +20.2 / +1.28 / +1.41 \\
first10 & 0.88 & 0.89 & 0.47 & 0.79 & 0.35 & +18.7 / +1.09 / +1.29 \\
first10\_x3 & 0.98 & 0.98 & 0.47 & 0.90 & 0.34 & +17.4 / +1.08 / +1.15 \\
\bottomrule
\end{tabular}
\end{table}

\begin{table}[htbp]
\centering
\scriptsize
\caption{27B: as Table~\ref{tab:app-meanabl12}.}
\label{tab:app-meanabl27}
\begin{tabular}{@{}lrrrrrr@{}}
\toprule
condition & own step only & agree & \begin{tabular}[b]{@{}c@{}}first\\{}step only\end{tabular} & \begin{tabular}[b]{@{}c@{}}conflict,\\{}C larger\end{tabular} & \begin{tabular}[b]{@{}c@{}}conflict,\\{}M larger\end{tabular} & \begin{tabular}[b]{@{}c@{}}regression:\\{}own\_yes / d1 / d2\end{tabular} \\
\midrule
clean & 0.99 & 0.91 & 0.50 & 0.98 & 0.23 & +35.4 / +1.74 / +2.15 \\
own5 & 0.99 & 0.91 & 0.50 & 0.98 & 0.23 & +39.1 / +1.85 / +2.40 \\
own10 & 0.99 & 0.90 & 0.49 & 0.97 & 0.23 & +37.4 / +1.86 / +2.34 \\
own20 & 0.97 & 0.89 & 0.51 & 0.92 & 0.26 & +35.3 / +1.81 / +2.44 \\
random & 0.99 & 0.91 & 0.50 & 0.98 & 0.23 & +35.3 / +1.74 / +2.14 \\
first10 & 0.99 & 0.91 & 0.49 & 0.97 & 0.23 & +42.5 / +2.04 / +2.59 \\
first10\_x3 & 0.99 & 0.90 & 0.50 & 0.97 & 0.23 & +28.9 / +1.48 / +1.87 \\
\bottomrule
\end{tabular}
\end{table}

\begin{table}[htbp]
\centering
\scriptsize
\caption{12B, removal set: accuracy and regression coefficients under twin-mean removal, swap and controls.}
\label{tab:app-removal12}
\begin{tabular}{@{}lrrrrrrrr@{}}
\toprule
condition & own step only & agree & \begin{tabular}[b]{@{}c@{}}conflict,\\{}C larger\end{tabular} & \begin{tabular}[b]{@{}c@{}}conflict,\\{}M larger\end{tabular} & shortcut & truth & d1 & d2 \\
\midrule
clean & 0.95 & 0.97 & 0.86 & 0.32 & +20.3 & +1.5 & +1.06 & +0.81 \\
twinmean\_top20 & 0.89 & 0.89 & 0.71 & 0.46 & +6.2 & +0.4 & +1.11 & +0.19 \\
twinmean\_rand20 & 0.96 & 0.98 & 0.87 & 0.29 & +28.0 & +2.1 & +1.01 & +1.02 \\
twinmean\_top10 & 0.93 & 0.96 & 0.82 & 0.39 & +11.0 & +0.8 & +1.08 & +0.41 \\
twinmean\_rand10 & 0.96 & 0.97 & 0.86 & 0.31 & +22.0 & +1.7 & +1.05 & +0.87 \\
swap\_top20 & 0.02 & 0.28 & 0.01 & 0.93 & -10.8 & -1.4 & +1.23 & -0.52 \\
swap\_top10 & 0.29 & 0.57 & 0.18 & 0.76 & -1.2 & -0.5 & +1.21 & -0.19 \\
linear prediction & 0.59 & 0.61 & 0.54 & 0.47 & 0 &  &  &  \\
\bottomrule
\end{tabular}
\end{table}

\begin{table}[htbp]
\centering
\scriptsize
\caption{27B: as Table~\ref{tab:app-removal12}.}
\label{tab:app-removal27}
\begin{tabular}{@{}lrrrrrrrr@{}}
\toprule
condition & own step only & agree & \begin{tabular}[b]{@{}c@{}}conflict,\\{}C larger\end{tabular} & \begin{tabular}[b]{@{}c@{}}conflict,\\{}M larger\end{tabular} & shortcut & truth & d1 & d2 \\
\midrule
clean & 0.98 & 0.86 & 0.98 & 0.21 & +33.3 & +6.1 & -0.16 & +0.59 \\
twinmean\_top20 & 0.98 & 0.86 & 0.96 & 0.29 & +18.6 & +4.4 & +0.10 & +0.60 \\
twinmean\_rand20 & 0.98 & 0.86 & 0.98 & 0.21 & +33.8 & +6.1 & -0.12 & +0.64 \\
twinmean\_top10 & 0.98 & 0.86 & 0.97 & 0.27 & +21.1 & +4.7 & +0.05 & +0.53 \\
twinmean\_rand10 & 0.98 & 0.86 & 0.98 & 0.21 & +33.4 & +6.2 & -0.16 & +0.58 \\
swap\_top20 & 0.63 & 0.73 & 0.51 & 0.58 & +1.0 & +1.5 & +0.70 & +0.27 \\
swap\_top10 & 0.91 & 0.83 & 0.79 & 0.43 & +7.2 & +2.7 & +0.47 & +0.43 \\
linear prediction & 0.76 & 0.65 & 0.58 & 0.51 & 0 &  &  &  \\
\bottomrule
\end{tabular}
\end{table}

\begin{table}[htbp]
\centering
\scriptsize
\caption{12B: where the removed shortcut comes back (change in the own\_yes coefficient by group, exact attribution on 600 questions).}
\label{tab:app-repair12}
\begin{tabular}{@{}lrrrrrrr@{}}
\toprule
removed & \begin{tabular}[b]{@{}c@{}}shortcut clean\\{}→ removed\end{tabular} & \begin{tabular}[b]{@{}c@{}}removed\\{}components\end{tabular} & \begin{tabular}[b]{@{}c@{}}other heads,\\{}same layers\end{tabular} & \begin{tabular}[b]{@{}c@{}}later\\{}components\end{tabular} & \begin{tabular}[b]{@{}c@{}}earlier\\{}(should be 0)\end{tabular} & final norm & \begin{tabular}[b]{@{}c@{}}final-norm scale\\{}removed ÷ clean\end{tabular} \\
\midrule
top20 & +18.1 → +4.3 & -46.3 & +15.7 & +15.5 & -0.00 & +1.3 & 1.119 \\
top10 & +18.1 → +9.2 & -34.2 & +9.3 & +14.2 & -0.00 & +1.9 & 1.143 \\
\bottomrule
\end{tabular}
\end{table}

\begin{table}[htbp]
\centering
\scriptsize
\caption{27B: as Table~\ref{tab:app-repair12}.}
\label{tab:app-repair27}
\begin{tabular}{@{}lrrrrrrr@{}}
\toprule
removed & \begin{tabular}[b]{@{}c@{}}shortcut clean\\{}→ removed\end{tabular} & \begin{tabular}[b]{@{}c@{}}removed\\{}components\end{tabular} & \begin{tabular}[b]{@{}c@{}}other heads,\\{}same layers\end{tabular} & \begin{tabular}[b]{@{}c@{}}later\\{}components\end{tabular} & \begin{tabular}[b]{@{}c@{}}earlier\\{}(should be 0)\end{tabular} & final norm & \begin{tabular}[b]{@{}c@{}}final-norm scale\\{}removed ÷ clean\end{tabular} \\
\midrule
top20 & +36.9 → +21.9 & -52.6 & +11.4 & +22.8 & +0.00 & +3.4 & 1.159 \\
top10 & +36.9 → +24.4 & -37.1 & +4.5 & +17.4 & +0.00 & +2.6 & 1.098 \\
\bottomrule
\end{tabular}
\end{table}

\begin{table}[htbp]
\centering
\scriptsize
\caption{Order-and-voice control (12B): rule coefficients from the pooled regression.}
\label{tab:app-ov-coef}
\begin{tabular}{@{}lr@{}}
\toprule
feature & \begin{tabular}[b]{@{}c@{}}coefficient\\{}[95\% CI]\end{tabular} \\
\midrule
own\_yes & +8.03 [+6.85, +9.20] \\
word\_yes & +11.19 [+10.66, +11.72] \\
last\_yes & -1.32 [-1.84, -0.78] \\
\bottomrule
\end{tabular}
\end{table}

\begin{table}[htbp]
\centering
\scriptsize
\caption{Order-and-voice control: accuracy by kind of question and rendering (MC = C's sentence last; CM = first; subj / obj = C as subject / object of its sentence).}
\label{tab:app-ov-acc}
\begin{tabular}{@{}lr@{}}
\toprule
difference & \begin{tabular}[b]{@{}c@{}}estimate\\{}[95\% CI]\end{tabular} \\
\midrule
own − word & -3.15 [-4.52, -1.88] \\
own − last & +9.36 [+8.09, +10.60] \\
word − last & +12.51 [+11.60, +13.33] \\
\bottomrule
\end{tabular}
\end{table}

\begin{table}[htbp]
\centering
\scriptsize
\caption{Order-and-voice control: share of answers C by own\_yes, word\_yes and last\_yes.}
\label{tab:app-ov-split}
\begin{tabular}{@{}lrrrr@{}}
\toprule
kind & MC-subj & CM-subj & MC-obj & CM-obj \\
\midrule
own step only & 0.97 & 0.97 & 0.41 & 0.47 \\
agree & 0.95 & 0.91 & 0.59 & 0.54 \\
first step only & 0.49 & 0.52 & 0.45 & 0.44 \\
conflict, C larger & 0.86 & 0.88 & 0.35 & 0.62 \\
conflict, M larger & 0.31 & 0.18 & 0.54 & 0.38 \\
\bottomrule
\end{tabular}
\end{table}

\begin{table}[htbp]
\centering
\scriptsize
\caption{Order-and-voice control: decisive questions by rendering.}
\label{tab:app-ov-dec}
\begin{tabular}{@{}lrrrrr@{}}
\toprule
rendering & own & word & last & n & answers C \\
\midrule
CM-obj & 0 & 0 & 0 & 140 & 0.49 \\
CM-obj & 0 & 0 & 1 & 160 & 0.38 \\
CM-obj & 0 & 1 & 0 & 288 & 0.81 \\
CM-obj & 0 & 1 & 1 & 456 & 0.61 \\
CM-obj & 1 & 0 & 0 & 590 & 0.87 \\
CM-obj & 1 & 0 & 1 & 166 & 0.69 \\
CM-subj & 0 & 0 & 0 & 428 & 0.02 \\
CM-subj & 0 & 0 & 1 & 616 & 0.04 \\
CM-subj & 1 & 1 & 0 & 590 & 0.75 \\
CM-subj & 1 & 1 & 1 & 166 & 0.85 \\
MC-obj & 0 & 0 & 0 & 300 & 0.03 \\
MC-obj & 0 & 1 & 1 & 744 & 0.33 \\
MC-obj & 1 & 0 & 0 & 756 & 0.25 \\
MC-subj & 0 & 0 & 0 & 1044 & 0.01 \\
MC-subj & 1 & 1 & 1 & 756 & 0.68 \\
\bottomrule
\end{tabular}
\end{table}

\begin{table}[htbp]
\centering
\scriptsize
\caption{Cross-family check (Qwen3-8B, Llama 3.1 8B): the preregistered criteria. Probe and model on the same held-out decisive questions, Wilson 95\% intervals.}
\label{tab:app-families-full}
\begin{tabular}{@{}lrrrrrrr@{}}
\toprule
model & layers & \begin{tabular}[b]{@{}c@{}}decisive\\{}accuracy\end{tabular} & \begin{tabular}[b]{@{}c@{}}R² d1 / d2 at the\\{}final position\end{tabular} & \begin{tabular}[b]{@{}c@{}}probe (best\\{}layer)\end{tabular} & \begin{tabular}[b]{@{}c@{}}model, same\\{}questions\end{tabular} & n & verdict \\
\midrule
Qwen3-8B & 36 & 0.36 & 0.78 / 0.75 & 0.86 [0.78, 0.91] (layer 24) & 0.25 [0.18, 0.34] & 104 & \textbf{replicates} \\
Llama 3.1 8B & 32 & 0.18 & 0.81 / 0.83 & 0.91 [0.84, 0.95] (layer 18) & 0.14 [0.09, 0.22] & 104 & \textbf{replicates} \\
\bottomrule
\end{tabular}
\end{table}

\begin{table}[htbp]
\centering
\scriptsize
\caption{Cross-family check: clean accuracy by kind of question.}
\label{tab:app-families-acc}
\begin{tabular}{@{}lrrrrr@{}}
\toprule
model & own step only & agree & \begin{tabular}[b]{@{}c@{}}first\\{}step only\end{tabular} & \begin{tabular}[b]{@{}c@{}}conflict,\\{}C larger\end{tabular} & \begin{tabular}[b]{@{}c@{}}conflict,\\{}M larger\end{tabular} \\
\midrule
Qwen3-8B & 0.70 & 0.77 & 0.54 & 0.69 & 0.36 \\
Llama 3.1 8B & 0.76 & 0.89 & 0.48 & 0.93 & 0.18 \\
\bottomrule
\end{tabular}
\end{table}

